\section{Masures: general framework}
We recollect here some well-known facts. Further details are available in the first two sections of~\cite{jep88refRou11}.
\subsection{Root generating system and Weyl groups}
A Kac--Moody matrix (or generalized Cartan matrix) is a square matrix $A=(a_{i,j})_{i,j\in I}$ indexed by a finite set $I$, with integral coefficients, and such that:
\begin{enumerate}[label=(\roman*)]
\item $\forall\ i\in I,\ a_{i,i}=2$;
\item $\forall\ (i,j)\in I^2,\ (i\neq j)\Rightarrow(a_{i,j}\leqslant0)$;
\item $\forall\ (i,j)\in I^2,\ (a_{i,j}=0)\Leftrightarrow(a_{j,i}=0$).
\end{enumerate}
A root generating system is a $5$-tuple $\jepPubScript{S}=(A,X,Y,(\alpha_i)_{i\in I},(\alpha_i^{\scriptscriptstyle\vee})_{i\in I})$ made of a Kac--Moody matrix $A$ indexed by the finite set $I$, of two dual free $\jepPubBbb{Z}$-modules $X$ and $Y$ of finite rank, and of a free family $(\alpha_i)_{i\in I}$ (respectively $(\alpha_i^{\scriptscriptstyle\vee})_{i\in I}$) of elements in $X$ (resp. $Y$) called simple roots (resp. simple coroots) that satisfy $a_{i,j}=\alpha_j(\alpha_i^{\scriptscriptstyle\vee})$ for all $i,j$ in $I$. Elements of $X$ (respectively of $Y$) are called characters (resp. cocharacters).

Fix such a root generating system $\jepPubScript{S}=(A,X,Y,(\alpha_i)_{i\in I},(\alpha_i^{\scriptscriptstyle\vee})_{i\in I})$ and set $\jepPubBbb{A}:=Y\otimes\jepPubBbb{R}$. Each element of $X$ induces a linear form on $\jepPubBbb{A}$, hence $X$ can be seen as a subset of the dual $\jepPubBbb{A}^*$. In particular, the $\alpha_i$'s (with $i\in I$) will be seen as linear forms on $\jepPubBbb{A}$. This allows us to define, for any $i\in I$, an involution $r_i$ of $\jepPubBbb{A}$ by setting $r_i(v):=v-\alpha_i(v)\alpha_i^{\scriptscriptstyle\vee}$ for any $v\in\jepPubBbb{A}$. Note that the points fixed by $r_i$ are exactly the elements of $\ker\alpha_i$. We define the Weyl group of $\jepPubScript{S}$ as the subgroup $W^v$ of $\mathrm{GL}(\jepPubBbb{A})$ generated by the finite set $\{r_i,\ i\in I\}$. The pair $(W^v,\{r_i,\ i\in I\})$ is a Coxeter system, hence we can consider the length $\ell(w)$ with respect to $\{r_i,\ i\in I\}$ of any element $w$ of $W^v$.

For any $x\in\jepPubBbb{A}$, we set $\underline{\alpha}(x)=(\alpha_i(x))_{i\in I}\in\jepPubBbb{R}^I$. Let $P^{\scriptscriptstyle\vee}:=\{v\in\jepPubBbb{A}\mid\underline{\alpha}(v)\in\jepPubBbb{Z}^I\}$ be the coweight-lattice, which is not a lattice when $\jepPubBbb{A}_{\mathrm{in}}:=\bigcap_{i\in I}\ker\alpha_i$ is non-zero, $Q^{\scriptscriptstyle\vee}:=\bigoplus_{i\in I}\jepPubBbb{Z}\alpha_i^{\scriptscriptstyle\vee}$ be the coroot-lattice and $Q^{\scriptscriptstyle\vee}_{\jepPubBbb{R}}=\bigoplus_{i\in I}\jepPubBbb{R}\alpha_i^{\scriptscriptstyle\vee}$. Furthermore setting $Q^{\scriptscriptstyle\vee}_+:=\bigoplus_{i\in I}\jepPubBbb{N}\alpha_i^{\scriptscriptstyle\vee}$, we can define a pre-order $\leqslant_{Q^{\scriptscriptstyle\vee}}$ on $\jepPubBbb{A}$ as follows: for any $x,y\in\jepPubBbb{A}$, we say that $x\leqslant_{Q^{\scriptscriptstyle\vee}}y$ iff $y-x\in Q^{\scriptscriptstyle\vee}_+$. We also set $Q^{\scriptscriptstyle\vee}_-:=-Q^{\scriptscriptstyle\vee}_+$ for future reference.

There is an action of the Weyl group $W^v$ on $\jepPubBbb{A}^*$ given by the following formula:
\[
\forall\ x\in\jepPubBbb{A},\ w\in W^v,\ \alpha\in\jepPubBbb{A}^*,\qquad(w\cdot\alpha)(x):=\alpha(w^{-1}\cdot x).
\]
Let $\Phi:=\{w\cdot\alpha_i\mid(w,i)\in W^v\times I\}$ be the set or real roots: then $\Phi$ is a subset of $Q:=\bigoplus_{i\in I}\jepPubBbb{Z}\alpha_i$. We will also use the set $\Delta:=\Phi\cup\Delta_{\mathrm{im}}^+\cup\Delta_{\mathrm{im}}^-\subset Q$ of all roots, as defined in~\cite{jep88refKac94}. Note that $\Delta$ is stable under the action of $W^v$. For any root $\alpha\in\Delta$, we set $\Lambda_\alpha:=\jepPubBbb{Z}$ if $\alpha\in\Phi$, and $\Lambda_\alpha:=\jepPubBbb{R}$ otherwise (i.e., if $\alpha\in\Delta_{\mathrm{im}}:=\Delta_{\mathrm{im}}^+\cup\Delta_{\mathrm{im}}^-$). For any pair $(\alpha,k)\in\Delta\times\jepPubBbb{Z}$, we set
\begin{equation}\label{jep88localHalfspaces}
D(\alpha,k):=\{t\in\jepPubBbb{A}\mid\alpha(t)+k\geqslant0\}\quad\text{and}\quad D^\circ(\alpha,k):=\{t\in\jepPubBbb{A}\mid\alpha(t)+k>0\}.
\end{equation}
For any root $\alpha\in\Delta$, we also set $D(\alpha,+\infty):=\jepPubBbb{A}$.

Finally, we let $W:=Q^{\scriptscriptstyle\vee}\rtimes W^v\subset\mathrm{GA}(\jepPubBbb{A})$ be the affine Weyl group of $\jepPubScript{S}$, where $\mathrm{GA}(\jepPubBbb{A})$ denotes the group of affine isomorphisms of $\jepPubBbb{A}$. Note that $W\subset P^{\scriptscriptstyle\vee}\rtimes W^v$ and that $\alpha(P^{\scriptscriptstyle\vee})$ is contained in $\jepPubBbb{Z}$ for any $\alpha\in Q$. Consequently, if $\tau$ is a translation of $\jepPubBbb{A}$ of vector $p\in P^{\scriptscriptstyle\vee}$, then for any $\alpha\in Q$, $\tau$ acts by permutations on the set $\{D(\alpha,k),\ k\in\jepPubBbb{Z}\}$. On the other hand, as $W^v$ stabilizes $\Delta$, any element of $W^v$ permutes the sets of the form $D(\alpha,k)$, where $\alpha$ runs over $\Delta$. Hence we have an action of $W$ on $\{D(\alpha,k),\ (\alpha,k)\in\Delta\times\jepPubBbb{Z}\}$.

% Exact publisher context draft physical6--7. Literal w dot x = 1 is retained as printed; no private-v2 substitution.
